\documentclass{article}
\usepackage{ifthen}
\usepackage{xparse}
\usepackage{etoolbox}
\usepackage{xspace}
\usepackage{amsmath}

\newcommand{\eg}{e.g.\xspace}
\newcommand{\ie}{i.e.,\xspace}
\newcommand{\vect}[1]{\mathbf{#1}}
\newcommand{\greet}[2][Hello]{#1, #2!}
\providecommand{\tightlist}{\setlength{\itemsep}{0pt}\setlength{\parskip}{0pt}}
\renewcommand{\vec}[1]{\boldsymbol{#1}}
\newcommand*{\note}[1]{\textit{(#1)}}
\def\half{\frac{1}{2}}
\def\pair#1#2{(#1, #2)}
\let\oldsection\section
\newcommand{\RR}{\mathbb{R}}
\NewDocumentCommand{\set}{m}{\{#1\}}
\newcommand{\draftmode}{true}

\newcounter{exercise}
\setcounter{exercise}{0}
\newenvironment{exercise}{\stepcounter{exercise}\par\noindent\textbf{Exercise \theexercise.} }{\par}
\newenvironment{important}[1][Important]{\begin{center}\textbf{#1:}}{\end{center}}

\begin{document}
\section{Macros}
Consider vectors, \eg $\vect{x} \in \RR^n$, \ie points. \greet{World} \greet[Goodbye]{Moon}
The half is $\half$ and the pair is $\pair{a}{b}$. Set $\set{1,2,3}$. \note{aside}
\tightlist

\begin{exercise}
First exercise.
\end{exercise}
\begin{exercise}
Second exercise; the counter is now \arabic{exercise} (roman: \roman{exercise}, alph: \alph{exercise}).
\end{exercise}
\addtocounter{exercise}{5}
Counter after adding five: \theexercise, value \the\value{exercise}.

\begin{important}
Read the docs.
\end{important}
\begin{important}[Warning]
Be careful.
\end{important}

\ifthenelse{\equal{\draftmode}{true}}{Draft mode is on.}{Final mode.}
\ifthenelse{\value{exercise} > 3}{Many exercises.}{Few exercises.}
\ifdefined\RR Defined.\fi

\newif\ifshowsolutions
\showsolutionstrue
\ifshowsolutions Solutions shown.\else Solutions hidden.\fi

\makeatletter
\renewcommand\@seccntformat[1]{\csname the#1\endcsname.\quad}
\makeatother

\newlength{\mylen}
\setlength{\mylen}{2cm}
\hspace{\mylen}Indented. \hspace*{1cm}More.
\vspace{1em}
\vspace*{\fill}

Text\smallskip

Text\medskip

Text\bigskip

\noindent\rule{\textwidth}{0.4pt}
\begin{flushright}
Right aligned.
\end{flushright}
\begin{flushleft}
Left aligned.
\end{flushleft}
{\raggedleft Ragged left.\par}
{\centering Centered text.\par}
\end{document}
